% Layout: transition (过渡衔接页)
% 使用场景: 两章之间的逻辑桥梁，承上启下
%
% Slots:
%   {{TITLE}}     - 页标题（如 "从H1到H2：为何...？"）
%   {{SUMMARY}}   - 上一章结论概括（50-80字）
%   {{QUESTIONS}} - 引出的问题/下一步方向（itemize格式，2-3条）
%
% 结构: 先总结→再提问→引出下一章

\begin{frame}
  \frametitle{{{TITLE}}}
  \vskip0.2cm

  {{SUMMARY}}

  \vskip0.25cm

  {{QUESTIONS}}
\end{frame}

% --- SUMMARY 示例 ---
%   \alert{H1 已确认形态边界的先导地位}。沿这一线索，自然引出第二层问题：
%   既然形态多样性是丰富度变化的领先信号，那么\alert{是什么驱动了形态边界本身的变化}？
%
% --- QUESTIONS 示例 ---
%   \begin{itemize}\setlength\itemsep{0.4em}
%     \item \textbf{短期外源冲击 (H2)}\,——\,环境扰动是否在形态边界上留下结构性变点？
%     \item \textbf{长期内在约束 (H3)}\,——\,演化轨迹为何能够回归而非持续恶化？
%   \end{itemize}
